\documentclass[11pt,a4paper]{article}
\usepackage{fontspec,amsmath,amssymb,amsthm,microtype}
\setmainfont{lmroman10-regular.otf}[BoldFont=lmroman10-bold.otf,ItalicFont=lmroman10-italic.otf,BoldItalicFont=lmroman10-bolditalic.otf]
\setmonofont{lmmono10-regular.otf}[Scale=MatchLowercase]
\usepackage[margin=25mm]{geometry}
\usepackage[unicode,hidelinks]{hyperref}
\setlength{\emergencystretch}{3em}
\setlength{\parskip}{0.4em}
\setlength{\parindent}{0pt}
\allowdisplaybreaks[2]
\raggedbottom
\title{\texorpdfstring{Identify actual Gaussian IPLA with the shared ULA construction}{Identify actual Gaussian IPLA with the shared ULA construction}}
\author{LangevinLean mathematical exposition}
\date{}
\begin{document}\maketitle
\noindent Source module: \texttt{IPLA/\allowbreak{}ULAAdapter.\allowbreak{}lean}. The accompanying HTML pairs this mathematical exposition with the exact Lean statements and records the scope of the formalization.


This adapter translates the particle coordinates into shared ULA probability interfaces. It proves the needed Gaussian and law identities, while all Brownian construction and stationarity arguments remain in their shared modules.


\section{\texorpdfstring{Brownian increments satisfy the complete algorithm noise specification}{Brownian increments satisfy the complete algorithm noise specification}}


On the independent product of an initial coupling and a Brownian source, define \(\xi_{k,a}=(B_a((k+1)h)-B_a(kh))/\sqrt h\). Each Brownian coordinate has centered Gaussian increments of variance \(h\), so each normalized coordinate has law \(\mathcal N(0,1)\). Independent Brownian coordinates and disjoint increments give mutual independence across coordinates and times. The product construction makes the complete noise path independent of the initial coupling. The shared product-noise theorems prove these three assertions, which are assembled into \texttt{Gaussian\allowbreak{}Innovations}.



The local instance \texttt{algorithm\allowbreak{}Product\allowbreak{}Probability} follows because a coupling of probability measures has total mass one and the independent source law has mass one; their product therefore has total mass one. This justifies the probability-space hypotheses of the innovation and L\ensuremath{{}^{2}} law theorems.




\textbf{Correspondence with the displayed declarations.}


\begin{itemize}

\item \texttt{product\allowbreak{}Gaussian\allowbreak{}Innovations}: The concrete product realization satisfies every distributional field of the actual IPLA algorithm. Standard Gaussian laws and all independence claims are conclusions from the independent Brownian coordinate source.
\end{itemize}


{\raggedright\noindent Lean declarations: \texttt{product\allowbreak{}Gaussian\allowbreak{}Innovations}, \texttt{algorithm\allowbreak{}Product\allowbreak{}Probability}.\par}

\section{\texorpdfstring{Transfer the product realization to any implementation}{Transfer the product realization to any implementation}}


The Brownian product's initial random variable has the prescribed law \(P\circ Z_0^{-1}\), and its innovations satisfy the Gaussian structure just proved. Implementation independence therefore identifies its IPLA law with any supplied Gaussian implementation having the same initial law. The shared product-chain law identity then gives

\[\rho_k=\operatorname{productOutputLaw}(V_N,\gamma/N,\mu_0,k).\]

Finite coupling cost and the target second moment are used here to construct the product initial state in \(L^2\). Neither stationarity nor a reference diffusion is assumed for this law identity.




Canonical Brownian Euler increments are proved standard Gaussian in every coordinate, mutually independent across times and coordinates, and independent of the coupled initialization. These are exactly the three fields of \texttt{Gaussian\allowbreak{}Innovations}. Their initial law agrees with the supplied initialization law on the product extension, so canonical-law uniqueness identifies the resulting algorithm output laws. The final Markov-step identity then shows that the actual supplied chain iterates the same ULA kernel at step \(\gamma/N\).




\textbf{Correspondence with the displayed declarations.}


\begin{itemize}

\item \texttt{laws\_\allowbreak{}eq\_\allowbreak{}product\allowbreak{}Brownian\allowbreak{}Output}: The law produced by the Brownian product construction is exactly the law of any supplied Gaussian IPLA implementation with the same initialization. This bridge has no stationarity or reference-process hypothesis.
\end{itemize}


{\raggedright\noindent Lean declarations: \texttt{laws\_\allowbreak{}eq\_\allowbreak{}product\allowbreak{}Brownian\allowbreak{}Output}.\par}

\section{\texorpdfstring{The actual chain iterates the ULA Markov transition}{The actual chain iterates the ULA Markov transition}}


The innovation hypotheses imply independent vector-valued Gaussian increments by grouping the scalar coordinate family. Applying the shared Hilbert-space law theorem to \(V_N\), step \(\gamma/N\), and smoothness \(NL\) gives

\[\rho_k=\operatorname{ulaLawStep}(V_N,\gamma/N)^{[k]}(\mu_0).\]

The continuous \(L^2\) representatives are measurable; the initial state is independent of the full noise path by hypothesis. This exact transition identity transfers both convergence and finiteness statements from ULA to the actual IPLA laws.




\textbf{Correspondence with the displayed declarations.}


\begin{itemize}

\item \texttt{laws\_\allowbreak{}eq\_\allowbreak{}ula\allowbreak{}Law\allowbreak{}Step}: The actual IPLA law is precisely ULA for the rescaled lifted potential. All vector Gaussian and independence facts follow from the displayed scalar innovation specification of the algorithm.
\end{itemize}


{\raggedright\noindent Lean declarations: \texttt{laws\_\allowbreak{}eq\_\allowbreak{}ula\allowbreak{}Law\allowbreak{}Step}.\par}

\end{document}
